% ============================================================
% Sección 7: Chain-of-Thought y ReAct
% ============================================================
\section{Chain-of-Thought y ReAct}

\begin{frame}{Chain-of-Thought (CoT)}
  \definicion{Chain-of-Thought}{
    Técnica de \emph{prompting} donde se pide (o se induce) al modelo a
    escribir su \textbf{razonamiento paso a paso} antes de dar la respuesta
    final, en lugar de responder directo.
  }
  \vfill
  \begin{itemize}
    \item Simple de aplicar: basta agregar ``pensemos paso a paso''\, al prompt.
    \item Mejora notablemente el desempeño en tareas de \textbf{lógica,
      matemáticas y problemas de varios pasos}.
    \item Es la base conceptual de los modelos ``que razonan''\ (sección anterior).
  \end{itemize}
\end{frame}

\begin{frame}{CoT en un diagrama}
  \centering
  \begin{tikzpicture}[node distance=0.5cm]
    \node[cajaAgente] (p) {Pregunta};
    \node[cajaLLM, below=of p, minimum width=6cm, minimum height=1.6cm] (r) {
      \begin{minipage}{5.4cm}\scriptsize
        \textbf{Razonamiento:} ``Primero calculo X, luego uso ese valor
        para obtener Y, y con eso...''
      \end{minipage}};
    \node[cajaAccento, below=of r] (a) {Respuesta final};
    \draw[flecha] (p) -- (r);
    \draw[flecha] (r) -- (a);
  \end{tikzpicture}
  \vfill
  \footnotesize CoT es puro \emph{razonamiento en texto} — todavía no
  interactúa con el mundo exterior.
\end{frame}

\begin{frame}{La limitación de CoT}
  \begin{itemize}
    \item CoT razona muy bien \textbf{sobre lo que ya sabe}.
    \item Pero no puede verificar hechos, hacer cálculos exactos grandes,
      ni consultar información nueva — solo ``piensa en voz alta''.
    \item ¿Qué pasa si el razonamiento necesita un dato externo
      (una búsqueda, un cálculo, un archivo)?
  \end{itemize}
  \vfill
  \centering
  \Large $\Rightarrow$ ReAct
\end{frame}

\begin{frame}{ReAct: Reasoning + Acting}
  \definicion{ReAct (Yao et al., 2022)}{
    Intercalar pasos de \textbf{razonamiento} (Thought) con pasos de
    \textbf{acción} (Action) sobre herramientas externas, observando el
    resultado (Observation) antes de continuar razonando.
  }
  \vfill
  \small
  \texttt{Thought} $\to$ \texttt{Action} $\to$ \texttt{Observation} $\to$
  \texttt{Thought} $\to$ \dots $\to$ \texttt{Respuesta final}
\end{frame}

\begin{frame}{El loop ReAct}
  \centering
  \begin{tikzpicture}[node distance=1.1cm]
    \node[cajaLLM] (t) {\textbf{Thought}\\¿qué necesito hacer?};
    \node[cajaHerramienta, right=1.8cm of t] (a) {\textbf{Action}\\llamar herramienta};
    \node[cajaEntorno, below=of a] (o) {\textbf{Observation}\\resultado obtenido};
    \node[cajaLLM, below=of t] (t2) {\textbf{Thought}\\¿ya tengo la respuesta?};

    \draw[flecha] (t) -- (a);
    \draw[flecha] (a) -- (o);
    \draw[flecha] (o) -- (t2);
    \draw[flecha] (t2.west) to[out=180,in=180,looseness=1.3] node[left, etiquetaFlecha, align=center]{repetir\\si falta info} (t.west);
    \node[cajaAccento, below=of t2] (fin) {Respuesta final};
    \draw[flecha] (t2) -- node[right, etiquetaFlecha]{si ya alcanza} (fin);
  \end{tikzpicture}
\end{frame}

\begin{frame}{Ejemplo de traza ReAct}
  \footnotesize
  \begin{block}{Pregunta}
    ``¿Qué día de la semana cae el 25 de diciembre de este año, y a cuánto
    equivale en horas si faltan 40 días?''
  \end{block}
  \begin{itemize}
    \item \texttt{Thought}: necesito saber la fecha de hoy.
    \item \texttt{Action}: \texttt{fecha\_actual()}
    \item \texttt{Observation}: \texttt{"2026-07-02"}
    \item \texttt{Thought}: ahora calculo los días restantes y la conversión.
    \item \texttt{Action}: \texttt{calculadora("40 * 24")}
    \item \texttt{Observation}: \texttt{960}
    \item \texttt{Thought}: ya puedo responder.
    \item \textbf{Respuesta final}: ``Faltan 960 horas; el 25 de diciembre es viernes.''
  \end{itemize}
\end{frame}
